% GENERATED FILE. Edit canonical lessons, then run npm run generate:books.
% canonical-source-hash: fnv1a64:b42aa24ef0f48b00
% canonical-lessons: RU-C431-konditsioner, RU-C431-obogrevatel, RU-C431-skvoznyak, RU-C431-provetrivat, RU-C431-sushit

\chapter{Too hot, too cold}
\label{ch:ru-too-hot-too-cold}
\hlchaptermodality{431}

\begin{chapteropening}
\textbf{By the end of this chapter:} \emph{I can deal with heat and cold in a room: the air conditioner, the heater, a draught, airing the room and drying things.}

Complete the last lesson of chapter 431: I can deal with heat and cold in a room: the air conditioner, the heater, a draught, airing the room and drying things.
\end{chapteropening}

\section[\texorpdfstring{\ru{кондиционер} (konditsionér)}{кондиционер (konditsionér)}]{\ru{кондиционер} (konditsionér) — an air conditioner; a hair conditioner}

\label{lesson:RU-C431-konditsioner}



\begin{quote}
Before the new one: say the Russian for a chambermaid, then the Russian for to make up.
\end{quote}

\subsection*{You'll want to know: \ru{кондиционер}}
\textbf{\ru{кондиционер}} (\emph{konditsionér}) — \textquotedblleft{}an air conditioner; a hair conditioner\textquotedblright{}.

\begin{quote}
\ru{В} \ru{жаркие} \ru{дни} \ru{кондиционер} \ru{работает} \ru{весь} \ru{день}.
\end{quote}

\begin{quote}
\emph{V zhárkiye dni konditsionér rabótayet ves den.}
\end{quote}

\begin{quote}
On hot days the air conditioner runs all day.
\end{quote}

Masculine.

\begin{cousinweb}
\emph{\ru{кондиционер}} was borrowed from English \emph{air conditioner} (EtymDB), keeping the second word.
\end{cousinweb}

\begin{grammarlens}[title={What you've built}]
One more word for telling, explaining and describing at B1.
\end{grammarlens}

\subsection*{Your turn}
\begin{itemize}
\raggedright
  \item \emph{Say it:} \emph{konditsionér}
  \item \emph{Say it:} \emph{konditsionér}, once more
  \item \emph{Say it:} a sentence of your own with \emph{konditsionér}
  \item \emph{Recall:} say the Russian for a chambermaid, then the Russian for to make up, then say \emph{konditsionér} again
\end{itemize}

\begin{tcolorbox}[breakable,colback=teal!4,colframe=teal!35!black,title={Before you move on}]
What does \ru{кондиционер} mean? (\textquotedblleft{}An air conditioner; a hair conditioner\textquotedblright{}.) Say it once more.
\end{tcolorbox}
\section[\texorpdfstring{\ru{обогреватель} (obogrevátel)}{обогреватель (obogrevátel)}]{\ru{обогреватель} (obogrevátel) — a heater}

\label{lesson:RU-C431-obogrevatel}



\begin{quote}
Before the new one: say the Russian for to make up, then the Russian for an air conditioner.
\end{quote}

\subsection*{You'll want to know: \ru{обогреватель}}
\textbf{\ru{обогреватель}} (\emph{obogrevátel}) — \textquotedblleft{}a heater\textquotedblright{}.

\begin{quote}
\ru{Ночью} \ru{в} \ru{доме} \ru{было} \ru{холодно}, \ru{поэтому} \ru{мы} \ru{включили} \ru{обогреватель}.
\end{quote}

\begin{quote}
\emph{Nóchyu v dóme býlo khólodno, poétomu my vklyuchíli obogrevátel.}
\end{quote}

\begin{quote}
At night the house was cold, so we switched on the heater.
\end{quote}

Masculine, though it ends in -\ru{ь}.

\begin{cousinweb}
\emph{\ru{обогреватель}} is built on \emph{\ru{обогревать}} 'to warm up' with \emph{-\ru{тель}}, the thing that does it. \emph{\ru{Греть}} is inherited from Proto-Slavic \emph{grěti} 'to heat, to warm' (EtymDB).
\end{cousinweb}

\begin{grammarlens}[title={What you've built}]
One more word for telling, explaining and describing at B1.
\end{grammarlens}

\subsection*{Your turn}
\begin{itemize}
\raggedright
  \item \emph{Say it:} \emph{obogrevátel}
  \item \emph{Say it:} \emph{obogrevátel}, once more
  \item \emph{Say it:} a sentence of your own with \emph{obogrevátel}
  \item \emph{Recall:} say the Russian for to make up, then the Russian for an air conditioner, then say \emph{obogrevátel} again
\end{itemize}

\begin{tcolorbox}[breakable,colback=teal!4,colframe=teal!35!black,title={Before you move on}]
What does \ru{обогреватель} mean? (\textquotedblleft{}A heater\textquotedblright{}.) Say it once more.
\end{tcolorbox}
\section[\texorpdfstring{\ru{сквозняк} (skvoznyák)}{сквозняк (skvoznyák)}]{\ru{сквозняк} (skvoznyák) — a draught (of air)}

\label{lesson:RU-C431-skvoznyak}



\begin{quote}
Before the new one: say the Russian for an air conditioner, then the Russian for a heater.
\end{quote}

\subsection*{You'll want to know: \ru{сквозняк}}
\textbf{\ru{сквозняк}} (\emph{skvoznyák}) — \textquotedblleft{}a draught (of air)\textquotedblright{}.

\begin{quote}
\ru{Из} \ru{окна} \ru{дует}, \ru{в} \ru{комнате} \ru{сквозняк}.
\end{quote}

\begin{quote}
\emph{Iz okná dúyet, v kómnate skvoznyák.}
\end{quote}

\begin{quote}
Air is coming in through the window; there is a draught in the room.
\end{quote}

Masculine. Built on \emph{\ru{сквозь}} — through.

\begin{cousinweb}
\emph{\ru{сквозняк}} is built on \emph{\ru{сквозь}} 'through' with \emph{-\ru{няк}}: air that blows right through a room. EtymDB lists the word without a deeper source.
\end{cousinweb}

\begin{grammarlens}[title={What you've built}]
One more word for telling, explaining and describing at B1.
\end{grammarlens}

\subsection*{Your turn}
\begin{itemize}
\raggedright
  \item \emph{Say it:} \emph{skvoznyák}
  \item \emph{Say it:} \emph{skvoznyák}, once more
  \item \emph{Say it:} a sentence of your own with \emph{skvoznyák}
  \item \emph{Recall:} say the Russian for an air conditioner, then the Russian for a heater, then say \emph{skvoznyák} again
\end{itemize}

\begin{tcolorbox}[breakable,colback=teal!4,colframe=teal!35!black,title={Before you move on}]
What does \ru{сквозняк} mean? (\textquotedblleft{}A draught (of air)\textquotedblright{}.) Say it once more.
\end{tcolorbox}
\section[\texorpdfstring{\ru{проветривать} (provétrivat)}{проветривать (provétrivat)}]{\ru{проветривать} (provétrivat) — to air (a room)}

\label{lesson:RU-C431-provetrivat}



\begin{quote}
Before the new one: say the Russian for a heater, then the Russian for a draught.
\end{quote}

\subsection*{You'll want to know: \ru{проветривать}}
\textbf{\ru{проветривать}} (\emph{provétrivat}) — \textquotedblleft{}to air (a room)\textquotedblright{}.

\begin{quote}
\ru{Каждое} \ru{утро} \ru{мама} \ru{проветривает} \ru{комнату}.
\end{quote}

\begin{quote}
\emph{Kázhdoye útro máma provétrivayet kómnatu.}
\end{quote}

\begin{quote}
Every morning Mum airs the room.
\end{quote}

Imperfective; perfective \emph{\ru{проветрить}}. Built on \emph{\ru{ветер}} — wind.

\begin{cousinweb}
\emph{\ru{проветривать}} is built from \emph{\ru{про}-} 'through' and \emph{\ru{ветр}-}, the root of \emph{\ru{ветер}} 'wind': to let the wind through. EtymDB does not list the verb.
\end{cousinweb}

\begin{grammarlens}[title={What you've built}]
One more word for telling, explaining and describing at B1.
\end{grammarlens}

\subsection*{Your turn}
\begin{itemize}
\raggedright
  \item \emph{Say it:} \emph{provétrivat}
  \item \emph{Say it:} \emph{provétrivat}, once more
  \item \emph{Say it:} a sentence of your own with \emph{provétrivat}
  \item \emph{Recall:} say the Russian for a heater, then the Russian for a draught, then say \emph{provétrivat} again
\end{itemize}

\begin{tcolorbox}[breakable,colback=teal!4,colframe=teal!35!black,title={Before you move on}]
What does \ru{проветривать} mean? (\textquotedblleft{}To air (a room)\textquotedblright{}.) Say it once more.
\end{tcolorbox}
\section[\texorpdfstring{\ru{сушить} (sushít)}{сушить (sushít)}]{\ru{сушить} (sushít) — to dry}

\label{lesson:RU-C431-sushit}



\begin{quote}
Before the new one: say the Russian for a draught, then the Russian for to air.
\end{quote}

\subsection*{You'll want to know: \ru{сушить}}
\textbf{\ru{сушить}} (\emph{sushít}) — \textquotedblleft{}to dry\textquotedblright{}.

\begin{quote}
\ru{Мы} \ru{сушим} \ru{мокрые} \ru{вещи} \ru{на} \ru{балконе}.
\end{quote}

\begin{quote}
\emph{My súshim mókryye véshchi na balkóne.}
\end{quote}

\begin{quote}
We dry the wet things on the balcony.
\end{quote}

Imperfective; perfective \emph{\ru{высушить}}. Present \emph{\ru{сушу}, \ru{сушишь}}. \emph{\ru{Сушить} \ru{волосы}} — to dry one's hair.

\begin{cousinweb}
\emph{\ru{сушить}} is inherited from Proto-Slavic \emph{sušiti} (EtymDB), the 'make dry' verb of \emph{\ru{сухой}} 'dry'.
\end{cousinweb}

\begin{grammarlens}[title={What you've built}]
One more word for telling, explaining and describing at B1.
\end{grammarlens}

\subsection*{Your turn}
\begin{itemize}
\raggedright
  \item \emph{Say it:} \emph{sushít}
  \item \emph{Say it:} \emph{sushít}, once more
  \item \emph{Say it:} a sentence of your own with \emph{sushít}
  \item \emph{Recall:} say the Russian for a draught, then the Russian for to air, then say \emph{sushít} again
\end{itemize}

\begin{tcolorbox}[breakable,colback=teal!4,colframe=teal!35!black,title={Before you move on}]
What does \ru{сушить} mean? (\textquotedblleft{}To dry\textquotedblright{}.) Say it once more.
\end{tcolorbox}
